
\chapter{Entrevista}
\label{ap:entrevista_requerimientos}

\noindent
\textbf{Fecha:} 08-11-2026\\
\textbf{Lugar:} Panadería San Miguel, Santa Cruz de la Sierra, Bolivia\\
\textbf{Entrevistada:} Máxima Peres\\
\textbf{Cargo o función:} Propietaria de la panaderia \\

\par\medskip

\noindent
\textbf{1. ¿Qué tipos de pan elaboran y qué cantidades de ingredientes utilizan?}

\par\medskip
Hacemos pan casero, marraqueta y francés. Los sábados también hacemos empanadas y rollos de queso con la misma masa. Para el casero usamos 6,5 kilos de harina, 400 gramos de azúcar, 100 de sal y 30 de levadura. Para la marraqueta usamos 8 kilos de harina, 80 de azúcar, 150 de sal y 30 de levadura. El francés lleva 7,5 kilos de harina y las mismas cantidades de los otros ingredientes que la marraqueta. Para la masa del sábado usamos 5 kilos de harina, 150 de azúcar, 100 de sal y 20 de levadura. La levadura que usamos es en polvo.

\par\medskip

\noindent
\textbf{2. ¿Cuáles son las cantidades mínimas y máximas que necesitan dosificar?}

\par\medskip
Normalmente usamos entre 5 y 8 kilos de harina, de 80 a 400 gramos de azúcar, de 100 a 150 de sal y de 20 a 30 de levadura. Esas son las cantidades que manejamos, pero sería bueno poder cambiarlas cuando necesitemos preparar más o menos masa.

\par\medskip

\noindent
\textbf{3. ¿Cómo realizan actualmente la dosificación de los ingredientes?}

\par\medskip
Primero ponemos el balde en la balanza y la dejamos en cero. Después el trabajador va trayendo la harina en varias cargas hasta completar el peso. Luego agrega el azúcar, la sal y la levadura. Esos tres ingredientes los medimos con un envase pequeño, más o menos al cálculo. Cuando está todo listo, llevamos el balde a la amasadora.

\par\medskip

\noindent
\textbf{4. ¿Quién realiza este trabajo y qué dificultades tiene?}

\par\medskip
Lo hace el trabajador que prepara la masa. Lo que más cuesta es estar trayendo la harina y pendiente de la balanza. También tiene que ir por el azúcar. La sal y la levadura se ponen rápido, pero como las medimos al cálculo, a veces puede variar un poco la cantidad.

\par\medskip

\noindent
\textbf{5. ¿Qué le gustaría que hiciera una máquina dosificadora?}

\par\medskip
Me gustaría que uno ponga el balde, elija las cantidades y que la máquina vaya echando los ingredientes sola. Así el trabajador ya no tendría que estar llevando y pesando todo a mano. Pero tiene que ser fácil de usar y que se pueda sacar el balde sin problema.

\par\medskip

\noindent
\textbf{6. ¿Preferiría ingresar las cantidades manualmente o seleccionar una receta guardada?}

\par\medskip
Sería mejor tener las recetas guardadas, porque casi siempre hacemos los mismos panes. Por ejemplo, elegir pan casero y que ya salgan las cantidades. Pero también quisiera poder cambiarlas cuando haga falta.

\par\medskip

\noindent
\textbf{7. ¿Qué información le gustaría ver durante la dosificación?}

\par\medskip
Que nos muestre qué ingrediente está poniendo y cuánto falta para terminar. También que avise cuando ya esté listo para retirar el balde. Y si falta algún ingrediente o hay un problema, que nos diga qué está pasando.

\par\medskip

\noindent
\textbf{8. ¿Cuánto tiempo considera aceptable para dosificar los cuatro ingredientes?}

\par\medskip
No sabría decirle un tiempo exacto. Lo importante es que no nos retrase el trabajo. Si la máquina hace todo sola y el trabajador puede ocuparse de otra cosa mientras tanto, ya sería una buena ayuda.

\par\medskip

\noindent
\textbf{9. ¿Qué diferencia en las cantidades de los ingredientes considera aceptable?}

\par\medskip
No tenemos una cantidad exacta de error permitida. La harina sí la pesamos, pero los otros ingredientes los ponemos al cálculo. Lo que quisiera es que la máquina mida bien.
\par\medskip

\noindent
\textbf{10. ¿Cómo almacenan los ingredientes y cómo le gustaría recargar la máquina?}

\par\medskip
La harina viene en bolsas de 25 kilos. El azúcar también lo tenemos en bolsas, y la sal y la levadura están en la mesa. Me gustaría que sea fácil llenar la máquina y que no tengamos que levantar mucho peso. Podríamos recargarla entre una preparación y otra cuando haga falta.

\par\medskip

\noindent
\textbf{11. ¿Dónde considera que podría instalarse la dosificadora?}

\par\medskip
Podría ser por donde ahora pesamos los ingredientes, pero primero tendríamos que ver cuánto espacio ocupa. La amasadora está más al fondo y necesitamos dejar libre el paso para trabajar. No quisiera que la máquina estorbe para movernos o sacar los baldes.

\par\medskip

\noindent
\textbf{12. ¿Cómo le gustaría realizar la limpieza de la máquina?}

\par\medskip
Que sea fácil de limpiar, porque con la harina siempre queda polvo y restos. Sería bueno poder sacar las partes donde quedan los ingredientes, limpiarlas y volverlas a poner sin mucha complicación.

\par\medskip


\noindent
\textbf{13. ¿Qué considera más importante para que la dosificadora sea útil en la panadería?}

\par\medskip
Lo principal es que mida bien los ingredientes y que nos ahorre trabajo. También que sea fácil de manejar, limpiar y llenar. Que no ocupe demasiado espacio y que no tengamos que estar arreglándola a cada rato. Si realmente nos facilita preparar la masa, sería una buena ayuda.

\par\medskip

\noindent
\textbf{14. Para cada receta, ¿cuántos gramos de más o de menos aceptaría en la cantidad de harina por lote?}

Para empezar, diría que hasta 30 gramos de más o de menos en cualquiera de las recetas. Por ejemplo, para el pan casero usamos seis kilos y medio. Si queda un poquito por encima o por debajo, dentro de esos 30 gramos, me parece bien.

%------------------------------------------------------------------------------
%------------------------------------------------------------------------------
\begin{center}
    \vspace{1.5cm}
    \hspace*{0pt}
    \textcolor{gray}{\rule{180pt}{0.1pt}}
    \vspace{-0.2cm}\\
    \MakeUppercase{\bf{ Maxima Peres }}\\[-0.2cm]
    \MakeUppercase{\bf{ Propietaria de la panadería artesanal san miguel}}\\ 
\end{center}